%%% Tento soubor obsahuje definice různých užitečných maker a prostředí %%%
%%% Další makra připisujte sem, ať nepřekáží v ostatních souborech.     %%%

\usepackage{etoolbox}
\def\Jazyk{cze}
\def\ProhlaseniAIText{Prohlášení k~využití AI}
\let\ifstringequal\ifdefstring

% xurl musí být před hyperref (načítá ho pdfx), aby šly lámat dlouhé cesty a URL
\usepackage{xurl}
\usepackage{needspace}
\usepackage[a-2u]{pdfx}     % výsledné PDF bude ve standardu PDF/A-2u


%% Přepneme na českou sazbu, fonty Latin Modern a kódování češtiny
\usepackage[czech]{babel}
\usepackage{lmodern}
\usepackage{microtype}
\usepackage[T1]{fontenc}
\usepackage{textcomp}
\usepackage[utf8]{inputenc}
\usepackage{csquotes}

\appto{\bibsetup}{\sloppy}
\makeatletter
\def\UrlBigBreaks{\do\:}
\def\UrlBreaks{\do\/\do\.\do\-\do\?\do\,\do\_\do\%\do\=\do\#\do\~\do\&}
\makeatother

%%% Nastavení pro použití samostatné bibliografické databáze.
\usepackage[
   backend=biber
  ,style=apa
  ,apamaxprtauth=99
  %,sortlocale=cs_CZ
  ,sorting=nyt
  ,maxbibnames=20
  ,bibencoding=UTF8
]{biblatex}
\addbibresource{literatura.bib}
% Oficiální biblatex-apa neobsahuje czech-apa.lbx; v adresáři tex/ je odvozený soubor.
% Bez mapování na *-apa.lbx se v citacích místo roku vypisují názvy polí (např. labelday).
\DeclareLanguageMapping{czech}{czech-apa}
% V českém textu spojka mezi autory: „a“, ne anglické „and“ (i u stylu APA v těle práce).
\DefineBibliographyStrings{czech}{and = {a}}
\DeclareDelimFormat{finalnamedelim}{%
  \ifnumgreater{\value{listcount}}{\value{liststop}-1}
    {\addspace\bibstring{and}\space}
    {\addcomma\space}%
}

%%% Další užitečné balíčky (jsou součástí běžných distribucí LaTeXu)
\usepackage{amsmath}        % rozšíření pro sazbu matematiky
\usepackage{amsfonts}       % matematické fonty
\usepackage{amssymb}          % \checkmark a další symboly
\usepackage{amsthm}         % sazba vět, definic apod.
\usepackage{bm}             % tučné symboly (příkaz \bm)
\usepackage{graphicx}       % vkládání obrázků
% Grafy z kap.~3: složka img/results/ (Overleaf stačí nahrát tex/), případně experiments/... při kompilaci z tex/.
\graphicspath{{img/results/}{../experiments/results/figures/}{img/}}
\usepackage{fancyvrb}       % vylepšené prostředí pro strojové písmo
\usepackage{fancyhdr}       % prostředí pohodlnější nastavení hlavy a paty stránek
\usepackage{icomma}         % inteligetní čárka v matematickém módu
\usepackage{dcolumn}        % lepší zarovnání sloupců v tabulkách
\usepackage{booktabs}       % lepší vodorovné linky v tabulkách
\makeatletter
\@ifpackageloaded{xcolor}{
   \@ifpackagewith{xcolor}{usenames}{}{\PassOptionsToPackage{usenames}{xcolor}}
  }{\usepackage[usenames]{xcolor}} % barevná sazba
\makeatother
\usepackage{multicol}       % práce s více sloupci na stránce
\usepackage{tikz}           % tvorba diagramů
\usetikzlibrary{positioning, decorations.pathreplacing}
\usepackage{caption}
\usepackage{enumitem}
\setlist[itemize]{noitemsep, topsep=0pt, partopsep=0pt}
\setlist[enumerate]{noitemsep, topsep=0pt, partopsep=0pt}
\setlist[description]{noitemsep, topsep=0pt, partopsep=0pt}

\usepackage{tocloft}
\setlength\cftparskip{0pt}
\setlength\cftbeforechapskip{1.5ex}
\setlength\cftfigindent{0pt}
\setlength\cfttabindent{0pt}
\setlength\cftbeforeloftitleskip{0pt}
\setlength\cftbeforelottitleskip{0pt}
\setlength\cftbeforetoctitleskip{0pt}
\renewcommand{\cftlottitlefont}{\Huge\bfseries\sffamily}
\renewcommand{\cftloftitlefont}{\Huge\bfseries\sffamily}
\renewcommand{\cfttoctitlefont}{\Huge\bfseries\sffamily}

% vyznaceni odstavcu
\parindent=0pt
\parskip=11pt

% zakaz vdov a sirotku - jednoradkovych pocatku ci koncu odstavcu na prechodu mezi strankami
\clubpenalty=1000
\widowpenalty=1000
\displaywidowpenalty=1000

% nastaveni radkovani
\renewcommand{\baselinestretch}{1.20}

% nastaveni pro nadpisy - tucne a bezpatkove
\usepackage{sectsty}
\allsectionsfont{\sffamily}

% nastavení hlavy a paty stránek
\fancyhf{}
\fancyhead[R]{\rightmark}
\fancyfoot[R]{\thepage}
% \fancyhead[RO,LE]{\rightmark}
% \fancyfoot[RO,LE]{\thepage}
\renewcommand{\footrulewidth}{.5pt}
\fancypagestyle{plain}{%
\fancyhf{} % clear all header and footer fields
\fancyfoot[R]{\thepage}
% \fancyfoot[RO,LE]{\thepage}
\renewcommand{\headrulewidth}{0pt}
\renewcommand{\footrulewidth}{0.5pt}}

% Tato makra přesvědčují mírně ošklivým trikem LaTeX, aby hlavičky kapitol
% sázel příčetněji a nevynechával nad nimi spoustu místa. Směle ignorujte.
\makeatletter
\def\@makechapterhead#1{
  {\parindent \z@ \raggedright \sffamily
   \Huge\bfseries \thechapter. #1
   \par\nobreak
   \vskip 20\p@
}}
\def\@makeschapterhead#1{
  {\parindent \z@ \raggedright \sffamily
   \Huge\bfseries #1
   \par\nobreak
   \vskip 20\p@
}}
\makeatother

% Trochu volnější nastavení dělení slov, než je default.
\lefthyphenmin=2
\righthyphenmin=2

% Vypnuto: černé značky přetékajících řádků (overfull hbox)
\overfullrule=0pt

%% Balíček hyperref, kterým jdou vyrábět klikací odkazy v PDF,
%% ale hlavně ho používáme k uložení metadat do PDF (včetně obsahu).
%% Většinu nastavítek přednastaví balíček pdfx.
\hypersetup{unicode}
\hypersetup{breaklinks=true}
\hypersetup{hidelinks}

%%% Prostředí pro sazbu kódu, případně vstupu/výstupu počítačových
%%% programů. (Vyžaduje balíček fancyvrb -- fancy verbatim.)

\DefineVerbatimEnvironment{code}{Verbatim}{fontsize=\small, frame=single}

